\begin{figure*}[t]
\centering
\resizebox{\textwidth}{!}{%
\begin{tikzpicture}[
  every node/.style={font=\footnotesize},
  root/.style={box, fill=navy!12, draw=navy, font=\footnotesize\bfseries},
  fam/.style={hot, align=center},
  leaf/.style={plain, align=center, font=\scriptsize},
  edge from parent/.style={draw=steel, thick, -{Stealth[length=2mm]}},
  level 1/.style={sibling distance=64mm, level distance=14mm},
  level 2/.style={sibling distance=22mm, level distance=16mm},
]
\node[root] {Graph learning as a\\propagation operator $\Prop(\Lap)$}
  child {node[fam] {Classical\\(local / diffusive)}
    child {node[leaf] {spatial MP\\GCN, GAT,\\GraphSAGE, GIN}}
    child {node[leaf] {spectral\\ChebNet,\\graph wavelets}}
    child {node[leaf] {diffusion\\GDC, APPNP\\$e^{-t\Lap}$}}
  }
  child {node[fam] {Quantum\\(ballistic / unitary)}
    child {node[leaf] {quantum walks\\CTQW $e^{-it\Lap}$,\\QWNN}}
    child {node[leaf] {QGNN / PQC\\equivariant QGC,\\quantum kernels}}
  }
  child {node[fam] {Quantum-inspired\\(runnable today)}
    child {node[leaf] {$|e^{-it\Lap}|^2$ as a\\classical layer\\(this tutorial)}}
  }
  child {node[fam] {Dynamic-graph\\wrappers}
    child {node[leaf] {TGN, EvolveGCN,\\DySAT (any\\operator above)}}
  };
\end{tikzpicture}}
\caption{A taxonomy of graph learning organized by the propagation operator
$\Prop(\Lap)$. Classical operators are local (spatial message passing) or diffusive
(spectral / heat); quantum operators are unitary and ballistic (continuous-time quantum
walks, parameterized quantum graph circuits); the quantum-inspired branch runs the
quantum-walk kernel on classical hardware; dynamic-graph wrappers compose with any
operator to handle a time-varying topology. This tutorial threads the classical,
quantum, and quantum-inspired branches onto one axis.}
\label{fig:taxonomy}
\end{figure*}
